\chapter{การโต้ตอบเชิงพื้นที่และการตรวจจับโครงร่างมือ 21 จุดร่วม}
\label{chap:ch03}

\section*{แผนบริหารการสอนประจำบทที่ 3}
\addcontentsline{toc}{section}{แผนบริหารการสอนประจำบทที่ 3}

\noindent\textbf{หัวข้อเนื้อหาประจำบท}
\begin{enumerate}
    \item สถาปัตยกรรมส่วนขยายการติดตามมือ (XR\_EXT\_hand\_tracking)
    \item ลำดับชั้นกระดูกและพิกัดข้อต่อมือ 21 จุดร่วมในมาตรฐานสากล
    \item คณิตศาสตร์การอนุมานท่าทางมือ (Gesture Heuristics: Pinch, Grasp, Point)
    \item การออกแบบการปฏิสัมพันธ์แบบไร้สัมผัสกับอุปกรณ์วิทยาศาสตร์เสมือนจริง
\end{enumerate}

\noindent\textbf{วัตถุประสงค์เชิงพฤติกรรม}
\begin{enumerate}
    \item อธิบายกลไกของส่วนขยาย XR\_EXT\_hand\_tracking ใน OpenXR ได้
    \item ระบุตำแหน่งและความสัมพันธ์ของข้อต่อมือ 21 จุดร่วมได้อย่างถูกต้อง
    \item เขียนขั้นตอนวิธีคำนวณระยะหยิบจับ (Pinch Distance) และการจับยึดวัตถุได้
    \item พัฒนาระบบโต้ตอบเพื่อปรับตั้งค่าเครื่องมือทดลองเสมือนจริงได้อย่างราบรื่น
\end{enumerate}

\newpage

\noindent\textbf{กิจกรรมการเรียนการสอน}
\begin{enumerate}
    \item การบรรยายเชิงทฤษฎีพร้อมสาธิตการตรวจจับข้อต่อมือแบบสด (Live 3D Hand Mesh)
    \item กิจกรรมฝึกเขียนโปรแกรมอ่านพิกัดปลายนิ้วชี้และนิ้วหัวแม่มือ
    \item การทดลองจำลองการปรับหมุนสเกลไมโครมิเตอร์เสมือนด้วยการจีบนิ้ว
    \item การอภิปรายเรื่องเทคนิค Zero-GC เพื่อรักษาอัตราเฟรมคงที่ที่ 60 FPS
\end{enumerate}

\noindent\textbf{สื่อการเรียนการสอน}
\begin{enumerate}
    \item สไลด์บรรยายอิเล็กทรอนิกส์และแผนภาพกายวิภาคข้อต่อมือ 21 จุดร่วม
    \item สคริปต์ตัวอย่างการใช้งาน OpenXR Hand Tracking และ MediaPipe Hands
    \item ชุดทดสอบอุปกรณ์เสมือนจริง (Virtual Vernier Caliper \& Optical Slider)
    \item วิดีโอสาธิตการจับถืออุปกรณ์เคมีและฟิสิกส์ในห้องแล็บเสมือนจริง
\end{enumerate}

\noindent\textbf{การประเมินผล}
\begin{enumerate}
    \item การประเมินผลงานโค้ดตรวจจับท่าทางมือจีบ (Pinch Gesture Recognition)
    \item การประเมินความถูกต้องในการเชื่อมต่อลำดับกระดูกมือเข้ากับวัตถุเสมือน
    \item การทดสอบข้อเขียนและแบบฝึกหัดท้ายบท
\end{enumerate}

\newpage

\section{ส่วนขยายการติดตามมือและโครงสร้างข้อต่อ}

การปฏิสัมพันธ์ในห้องปฏิบัติการเสมือนจริงที่เป็นธรรมชาติสูงสุดคือการใช้มือเปล่าของผู้เรียนโดยไม่ต้องพึ่งพาคอนโทรลเลอร์ มาตรฐาน OpenXR รองรับความสามารถนี้ผ่านส่วนขยายทางการ \texttt{XR\_EXT\_hand\_tracking} ซึ่งทำหน้าที่รายงานพิกัดและทิศทางการหมุนของข้อต่อมือในระดับมิลลิเมตร

\begin{figure}[htbp]
\centering
\begin{tikzpicture}[scale=1.1]
    % Wrist
    \fill[darkNavy] (0, 0) circle (4pt) node[below=4pt]{\textbf{ข้อมือ (Wrist)}};

    % Metacarpals & Joints
    % Thumb
    \draw[line width=1.5pt, amberGlow] (0,0) -- (-1.2, 0.8) -- (-2.0, 1.6) -- (-2.6, 2.4);
    \fill[amberGlow] (-1.2, 0.8) circle (3pt);
    \fill[amberGlow] (-2.0, 1.6) circle (3pt);
    \fill[amberGlow] (-2.6, 2.4) circle (3.5pt) node[left]{\textbf{ปลายนิ้วหัวแม่มือ}};

    % Index
    \draw[line width=1.5pt, openxrTeal] (0,0) -- (-0.8, 1.8) -- (-0.9, 3.0) -- (-1.0, 4.2) -- (-1.1, 5.2);
    \fill[openxrTeal] (-0.8, 1.8) circle (3pt);
    \fill[openxrTeal] (-0.9, 3.0) circle (3pt);
    \fill[openxrTeal] (-1.0, 4.2) circle (3pt);
    \fill[openxrTeal] (-1.1, 5.2) circle (3.5pt) node[above]{\textbf{ปลายนิ้วชี้}};

    % Middle
    \draw[line width=1.5pt, emeraldActive] (0,0) -- (0.0, 2.0) -- (0.0, 3.4) -- (0.0, 4.6) -- (0.0, 5.8);
    \fill[emeraldActive] (0.0, 2.0) circle (3pt);
    \fill[emeraldActive] (0.0, 3.4) circle (3pt);
    \fill[emeraldActive] (0.0, 4.6) circle (3pt);
    \fill[emeraldActive] (0.0, 5.8) circle (3.5pt) node[above]{\textbf{ปลายนิ้วกลาง}};

    % Ring
    \draw[line width=1.5pt, cyberPurple] (0,0) -- (0.8, 1.8) -- (0.9, 3.1) -- (1.0, 4.3) -- (1.1, 5.3);
    \fill[cyberPurple] (0.8, 1.8) circle (3pt);
    \fill[cyberPurple] (0.9, 3.1) circle (3pt);
    \fill[cyberPurple] (1.0, 4.3) circle (3pt);
    \fill[cyberPurple] (1.1, 5.3) circle (3.5pt) node[above]{\textbf{ปลายนิ้วนาง}};

    % Little
    \draw[line width=1.5pt, spatialBlue] (0,0) -- (1.5, 1.4) -- (1.8, 2.5) -- (2.0, 3.5) -- (2.2, 4.4);
    \fill[spatialBlue] (1.5, 1.4) circle (3pt);
    \fill[spatialBlue] (1.8, 2.5) circle (3pt);
    \fill[spatialBlue] (2.0, 3.5) circle (3pt);
    \fill[spatialBlue] (2.2, 4.4) circle (3.5pt) node[right]{\textbf{ปลายนิ้วก้อย}};

    % Pinch Line
    \draw[dashed, line width=1.4pt, red!80!black] (-2.6, 2.4) -- (-1.1, 5.2) node[midway, above left]{\textbf{ระยะจีบ $d_{\text{pinch}}$}};
\end{tikzpicture}
\caption{โครงสร้างลำดับชั้นกระดูกมือ 21 จุดร่วมและการวัดระยะทางจีบเพื่อการปฏิสัมพันธ์}
\label{fig:hand_skeleton}
\end{figure}

ตามข้อกำหนดของ OpenXR ข้อต่อมือหลัก 21 จุดร่วม (และส่วนขยายปลีกย่อยรวม 26 จุด) ประกอบด้วยลำดับข้อต่อสำคัญ 5 นิ้ว ได้แก่ นิ้วหัวแม่มือ (Thumb), นิ้วชี้ (Index), นิ้วกลาง (Middle), นิ้วนาง (Ring), และนิ้วก้อย (Little) โดยแต่ละนิ้วประกอบด้วยข้อมือส่วนโคน (Metacarpal), ข้อกลาง (Proximal/Intermediate), ข้อปลาย (Distal), และปลายนิ้ว (Tip) ดังแสดงโครงสร้างและดัชนีจุดสังเกตในตารางที่ 3.1

\begin{table}[htbp]
\centering
\caption{ดัชนีและลำดับชั้นข้อต่อมือ 21 จุดร่วมตามมาตรฐาน MediaPipe และ OpenXR}
\label{tab:hand_landmarks_21}
\begin{tabularx}{\textwidth}{cllY}
\toprule
\textbf{ดัชนี} & \textbf{ชื่อข้อต่อ (Joint Name)} & \textbf{ส่วนของมือ} & \textbf{หน้าที่และบทบาทในการจำลองเสมือน} \\
\midrule
0 & \texttt{WRIST} & ข้อมือ & จุดกำเนิดอ้างอิงของมือ (Hand Root Coordinate) \\
1--4 & \texttt{THUMB\_CMC} ถึง \texttt{TIP} & นิ้วหัวแม่มือ & จุดที่ 4 ใช้คำนวณการจีบนิ้วและการหนีบวัตถุ \\
5--8 & \texttt{INDEX\_FINGER\_MCP} ถึง \texttt{TIP} & นิ้วชี้ & จุดที่ 8 ใช้เป็นตัวชี้ลำแสงและจุดจับยึดหลัก \\
9--12 & \texttt{MIDDLE\_FINGER\_MCP} ถึง \texttt{TIP} & นิ้วกลาง & ร่วมคำนวณการกำมือและทิศทางเวกเตอร์ระนาบมือ \\
13--16 & \texttt{RING\_FINGER\_MCP} ถึง \texttt{TIP} & นิ้วนาง & ตรวจจับการพับงอนิ้วเพื่อระบุท่าทาง Grasp \\
17--20 & \texttt{PINKY\_MCP} ถึง \texttt{TIP} & นิ้วก้อย & ตรวจจับขอบนอกสุดของฝ่ามือและระนาบรองรับ \\
\bottomrule
\end{tabularx}
\end{table}

\section{คณิตศาสตร์การอนุมานท่าทางมือ}

ในการควบคุมเครื่องมือทดลองเสมือนจริง ท่าทางพื้นฐานที่สำคัญที่สุดคือการจีบนิ้ว (Pinch Gesture) เพื่อหยิบจับวัตถุขนาดเล็กหรือปรับหมุนปุ่มลูกบิด และการกำมือ (Grasp Gesture) เพื่อยกเคลื่อนย้ายอุปกรณ์ขนาดใหญ่

\subsection{การคำนวณระยะการจีบนิ้ว}

ระยะห่างแบบยูคลิดระหว่างปลายนิ้วหัวแม่มือ $\mathbf{p}_{\text{thumb}}$ (ดัชนี 4) และปลายนิ้วชี้ $\mathbf{p}_{\text{index}}$ (ดัชนี 8) คำนวณได้จาก
\begin{equation}
d_{\text{pinch}} = \|\mathbf{p}_{\text{thumb}} - \mathbf{p}_{\text{index}}\| = \sqrt{(x_4 - x_8)^2 + (y_4 - y_8)^2 + (z_4 - z_8)^2}
\end{equation}
สถานะการจีบนิ้ว $S_{\text{pinch}}$ ถูกกำหนดด้วยฟังก์ชันฮิสเทอรีซิส (Hysteresis Thresholding) เพื่อป้องกันการสั่นไหวของสถานะที่จุดขอบเขต
\begin{equation}
S_{\text{pinch}}(t) = 
\begin{cases} 
\text{TRUE} & \text{ถ้า } d_{\text{pinch}} \le d_{\text{enter}} \\
\text{FALSE} & \text{ถ้า } d_{\text{pinch}} \ge d_{\text{exit}} \\
S_{\text{pinch}}(t-\Delta t) & \text{ในกรณีอื่น}
\end{cases}
\end{equation}
โดยทั่วไปกำหนดค่าเกณฑ์ $d_{\text{enter}} = 0.025\text{ m}$ ($2.5\text{ cm}$) และ $d_{\text{exit}} = 0.035\text{ m}$ ($3.5\text{ cm}$)

เมื่อสถานะ $S_{\text{pinch}} = \text{TRUE}$ จุดตำแหน่งกึ่งกลางของการหยิบจับ (Pinch Centroid) $\mathbf{P}_{\text{pinch}}$ ซึ่งใช้เป็นจุดเชื่อมต่อ (Interaction Anchor) สำหรับวัตถุเสมือน คำนวณได้จาก
\begin{equation}
\mathbf{P}_{\text{pinch}} = \frac{\mathbf{p}_{\text{thumb}} + \mathbf{p}_{\text{index}}}{2} = \left(\frac{x_4 + x_8}{2}, \frac{y_4 + y_8}{2}, \frac{z_4 + z_8}{2}\right)
\end{equation}

\subsection{ปฏิบัติการจำลองเครื่องมือวัดละเอียด: เวอร์เนียร์คาลิปเปอร์เสมือนจริง}

หนึ่งในนวัตกรรมการสอนสะเต็มศึกษาผ่านระบบความเป็นจริงขยาย คือการฝึกทักษะปฏิบัติการวัดละเอียดทางวิทยาศาสตร์ด้วย \textbf{เวอร์เนียร์คาลิปเปอร์เสมือนจริง (Virtual Vernier Caliper)} ผู้เรียนสามารถใช้นิ้วหัวแม่มือและนิ้วชี้จีบเลื่อนปากวัด (Sliding Jaw) เพื่อประกบวัตถุทดลอง เช่น ทรงกระบอกโลหะหรือลูกแก้วแก้ว

หลักการทำงานของการจำลองเวอร์เนียร์คาลิปเปอร์ประกอบด้วย 2 องค์ประกอบหลัก:
\begin{enumerate}
    \item \textbf{การจับยึดและเลื่อนแกนเดียว (1D Constrained Sliding):} พิกัดการจีบของมือจะถูกฉายลงบนแกนของสเกลหลัก $X_{\text{caliper}}$ เท่านั้น ปากวัดจะไม่หลุดออกจากระนาบแกนวัด
    \item \textbf{การคำนวณค่าความละเอียดและการอ่านค่า (Vernier Scale Reading):} เมื่อปากวัดเคลื่อนที่ไปเป็นระยะทาง $L$ การอ่านค่าประกอบด้วยค่าสเกลหลัก ($L_{\text{main}}$) และค่าขีดสเกลเวอร์เนียร์ที่ตรงกับสเกลหลักพอดี ($n$)
    \begin{equation}
    L_{\text{read}} = L_{\text{main}} + (n \times \text{LC})
    \end{equation}
    โดยที่ $\text{LC}$ คือค่าความละเอียดต่ำสุด (Least Count) ของเครื่องมือวัด เช่น $0.05\text{ mm}$ หรือ $0.01\text{ cm}$
\end{enumerate}

\begin{workedexamplebox}{การคำนวณสถานะการหยิบและอ่านค่าเวอร์เนียร์คาลิปเปอร์เสมือน}
ผู้เรียนใช้มือขวาควบคุมเวอร์เนียร์คาลิปเปอร์เสมือนเพื่อวัดความหนาของแผ่นอะคริลิก ปลายนิ้วหัวแม่มืออยู่ที่ $\mathbf{p}_4 = (0.15, 1.00, -0.40)\text{ m}$ และปลายนิ้วชี้อยู่ที่ $\mathbf{p}_8 = (0.17, 1.01, -0.39)\text{ m}$
\begin{enumerate}
    \item ตรวจสอบว่าเกิดการจีบนิ้วจับปากวัดหรือไม่ (เกณฑ์ $d_{\text{enter}} = 0.030\text{ m}$) และระบุจุดกึ่งกลางการจับยึด
    \item เมื่อปากวัดถูกเลื่อนไปประกบแผ่นอะคริลิก ขีดศูนย์ของสเกลเวอร์เนียร์อยู่ระหว่าง $18\text{ mm}$ ถึง $19\text{ mm}$ บนสเกลหลัก และขีดที่ 7 บนสเกลเวอร์เนียร์ตรงกับขีดบนสเกลหลักพอดี โดยเวอร์เนียร์มีความละเอียด $\text{LC} = 0.05\text{ mm}$ จงหาค่าความหนาของแผ่นอะคริลิก
\end{enumerate}

\textbf{วิธีทำ:}
1. คำนวณระยะห่างระหว่างนิ้ว
\begin{align}
d_{\text{pinch}} &= \sqrt{(0.15 - 0.17)^2 + (1.00 - 1.01)^2 + (-0.40 - (-0.39))^2} \notag \\
&= \sqrt{(-0.02)^2 + (-0.01)^2 + (-0.01)^2} = \sqrt{0.0004 + 0.0001 + 0.0001} \notag \\
&= \sqrt{0.0006} \approx 0.0245\text{ m} = 2.45\text{ cm}
\end{align}
เนื่องจาก $d_{\text{pinch}} = 0.0245\text{ m} < 0.030\text{ m}$ ระบบจึงระบุว่าเกิดการจีบนิ้วเพื่อจับปากวัด และมีตำแหน่งกึ่งกลางการจับยึดที่
\begin{align}
\mathbf{P}_{\text{pinch}} &= \left(\frac{0.15 + 0.17}{2}, \frac{1.00 + 1.01}{2}, \frac{-0.40 + (-0.39)}{2}\right) \notag \\
&= (0.160, 1.005, -0.395)\text{ m}
\end{align}

2. คำนวณค่าความหนาของแผ่นอะคริลิกจากการอ่านค่าเวอร์เนียร์
\begin{equation}
L_{\text{read}} = L_{\text{main}} + (n \times \text{LC}) = 18\text{ mm} + (7 \times 0.05\text{ mm}) = 18 + 0.35 = 18.35\text{ mm}
\end{equation}
แผ่นอะคริลิกมีความหนาเท่ากับ $18.35\text{ mm}$ หรือ $1.835\text{ cm}$
\end{workedexamplebox}

\section{วิศวกรรมลูปแบบไร้ขยะหน่วยความจำ}

ในการแสดงผลเชิงพื้นที่แบบ 60 ถึง 90 FPS การจัดสรรหน่วยความจำใหม่ซ้ำซากในทุก ๆ เฟรมจะส่งผลให้ระบบ Garbage Collector (GC) ทำงานและหยุดการทำงานชั่วขณะ (Micro-Stutters) ซึ่งทำลายความต่อเนื่องของการทดลอง นักพัฒนาต้องใช้หลักการ Zero-GC Loop Engineering โดยการสร้างโครงสร้างเวกเตอร์และตัวแปรคำนวณไว้ภายนอกลูปการเรนเดอร์ ดังแสดงในตารางที่ 3.1

\begin{table}[htbp]
\centering
\caption{แนวทางปฏิบัติในการจัดการหน่วยความจำสำหรับลูปคำนวณการติดตามมือ}
\label{tab:zero_gc_rules}
\begin{tabularx}{\textwidth}{lYY}
\toprule
\textbf{การกระทำที่ควรหลีกเลี่ยง} & \textbf{ผลกระทบต่อระบบ} & \textbf{แนวทางแก้ไขที่ถูกต้อง (Zero-GC)} \\
\midrule
สร้าง \texttt{new Vector3()} ในลูป & เกิดขยะอ็อบเจกต์ 90 ตัว/วินาที & จองตัวแปรเวกเตอร์ระดับโกลบอลเพียงครั้งเดียว \\
ต่อสตริงแสดงผล HUD ซ้ำ ๆ & เพิ่มภาระ Heap Memory & ใช้ Fixed-length Buffers หรือ Text Node เดียว \\
จัดสรรอาร์เรย์พิกัดใหม่ทุกเฟรม & อัตราเฟรมแกว่งและกระตุก & เขียนทับค่าลงใน TypedArray (Float32Array) เดิม \\
\bottomrule
\end{tabularx}
\end{table}

\begin{aipromptbox}[การสร้างระบบตรวจจับท่าทางมือแบบ Zero-GC ด้วย AI]
\textbf{วัตถุประสงค์:} ออกแบบคลาสสำหรับตรวจจับท่าทางการหยิบจับ (Pinch \& Grab) จากพิกัดข้อต่อมือ 21 จุดร่วมที่มีประสิทธิภาพสูง ปราศจากการจัดสรรหน่วยความจำใหม่ในลูปการเรนเดอร์ (Zero Garbage Collection)\\
\textbf{โครงสร้าง CLEAR Prompt:}
\begin{itemize}[topsep=1pt, itemsep=1pt, leftmargin=1.2em]
    \item \textbf{Context (บริบท):} ระบบโต้ตอบกับเครื่องมือวิทยาศาสตร์ใน OpenXR บนอุปกรณ์ Standalone ที่มีทรัพยากรหน่วยความจำจำกัด
    \item \textbf{Logic (ตรรกะ):} ใช้กลไก Hysteresis Thresholding เพื่อป้องกันการสั่นไหว พร้อมคำนวณเวกเตอร์ระยะห่างแบบ In-place บน \texttt{Float32Array} หรือ \texttt{np.ndarray} ที่จองล่วงหน้า
    \item \textbf{Expectation (ผลลัพธ์ที่ต้องการ):} คลาส Python หรือ JavaScript สำหรับติดตามสถานะ \texttt{PINCH\_START}, \texttt{PINCHING}, \texttt{PINCH\_END} พร้อมส่งพิกัดกึ่งกลางการหยิบจับ
    \item \textbf{Action \& Restriction (ข้อกำหนด):} ห้ามสร้าง Object หรือ Array ใหม่ในเมท็อด \texttt{update()} อย่างเด็ดขาด และต้องรักษาอัตราการคำนวณที่ $90\text{ FPS}$ ($< 1\text{ ms}$ ต่อเฟรม)
\end{itemize}
\end{aipromptbox}

\begin{codebox}[การคำนวณระยะจีบนิ้วแบบ Zero-GC ในภาษา Python]
import numpy as np

class HandPinchDetector:
    def __init__(self, enter_threshold=0.025, exit_threshold=0.035):
        self.enter_th = enter_threshold
        self.exit_th = exit_threshold
        self.is_pinching = False
        # จองบัฟเฟอร์หน่วยความจำล่วงหน้าเพื่อหลีกเลี่ยง GC
        self._diff = np.zeros(3, dtype=np.float32)

    def update(self, thumb_tip_pos, index_tip_pos):
        # คำนวณแบบ In-place บนอาร์เรย์เดิม
        np.subtract(thumb_tip_pos, index_tip_pos, out=self._diff)
        dist = np.linalg.norm(self._diff)

        if not self.is_pinching and dist <= self.enter_th:
            self.is_pinching = True
        elif self.is_pinching and dist >= self.exit_th:
            self.is_pinching = False

        return self.is_pinching, dist
\end{codebox}

\clearpage

\section*{สรุปท้ายบทที่ 3}
\addcontentsline{toc}{section}{สรุปท้ายบทที่ 3}

การตรวจจับโครงร่างมือ 21 จุดร่วมตามมาตรฐาน OpenXR มอบมิติใหม่ของการเรียนรู้สะเต็มศึกษา ช่วยให้ผู้เรียนสามารถควบคุม ปรับแต่ง และหยิบจับอุปกรณ์วิทยาศาสตร์ได้อย่างเป็นธรรมชาติ การประยุกต์ใช้คณิตศาสตร์ของระยะห่างแบบฮิสเทอรีซิสร่วมกับหลักการ Zero-GC ช่วยสร้างประสบการณ์ที่แม่นยำและเสถียร พร้อมสำหรับการเชื่อมต่อเข้ากับเอนจินฟิสิกส์ในบทต่อไป

\clearpage

\section*{แบบฝึกหัดท้ายบทที่ 3}
\addcontentsline{toc}{section}{แบบฝึกหัดท้ายบทที่ 3}

\begin{enumerate}
    \item จงอธิบายบทบาทของฟังก์ชันฮิสเทอรีซิสในการตรวจจับท่าทางการจีบนิ้ว และวิเคราะห์ผลกระทบหากใช้ค่าเกณฑ์ตัดสินเพียงค่าเดียว
    \item เขียนสมการคำนวณจุดกึ่งกลาง (Centroid) ของฝ่ามือโดยอ้างอิงจากตำแหน่งข้อต่อโคนนิ้วทั้งห้าจุด
    \item เพราะเหตุใดวิศวกรรมลูปแบบไร้ขยะหน่วยความจำ (Zero-GC Loop Engineering) จึงมีความจำเป็นอย่างยิ่งในการพัฒนาแอปพลิเคชันเสมือนจริงบนอุปกรณ์เคลื่อนที่
    \item จงระบุดัชนีของจุดข้อต่อปลายนิ้วหัวแม่มือและปลายนิ้วชี้ตามมาตรฐาน MediaPipe Hands และอธิบายวิธีการใช้สองจุดนี้ในการกำหนดจุดกึ่งกลางปฏิสัมพันธ์ (Interaction Anchor)
    \item ในการจำลองเวอร์เนียร์คาลิปเปอร์เสมือนจริง หากขีดศูนย์ของสเกลเวอร์เนียร์อยู่ระหว่าง $24\text{ mm}$ และ $25\text{ mm}$ บนสเกลหลัก และขีดที่ 6 บนสเกลเวอร์เนียร์ตรงกับขีดบนสเกลหลัก โดยมีความละเอียด $\text{LC} = 0.05\text{ mm}$ จงคำนวณค่าความยาวที่วัดได้
\end{enumerate}
